% GENERATED FILE. Edit canonical lessons, then run npm run generate:books.
% canonical-source-hash: fnv1a64:818235d72262e153
% canonical-lessons: KA-C457-oni, KA-C457-haduhogu, KA-C457-anvesisu, KA-C457-tana, KA-C457-kangolisu

\chapter{Wandering About}
\label{ch:ka-wandering-about}
\hlchaptermodality{457}

\begin{chapteropening}
\textbf{By the end of this chapter:} \emph{I can say a lane, to pass through, to explore, a spot and to look splendid.}

Complete the last lesson of chapter 457: I can say a lane, to pass through, to explore, a spot and to look splendid.
\end{chapteropening}

\section[\texorpdfstring{\kn{ಓಣಿ} (ōṇi)}{ಓಣಿ (ōṇi)}]{\kn{ಓಣಿ} (ōṇi) — a lane; an alley}

\label{lesson:KA-C457-oni}



\begin{quote}
Before the new one: say the Kannada for a landmark, then the Kannada for to view.
\end{quote}

\subsection*{You'll want to know: \kn{ಓಣಿ}}
\textbf{\kn{ಓಣಿ}} — \emph{ōṇi} — \textquotedblleft{}a lane; an alley\textquotedblright{}.

\begin{quote}
\kn{ಹಳೆಯ} \kn{ಊರಿನ} \kn{ಓಣಿಗಳು} \kn{ತುಂಬಾ} \kn{ಕಿರಿದು}.
\end{quote}

\begin{quote}
\emph{haḷeya ūrina ōṇigaḷu tumbā kiridu.}
\end{quote}

\begin{quote}
The lanes of the old town are very narrow.
\end{quote}

Plural: \kn{ಓಣಿಗಳು}, \emph{ōṇigaḷu}.

A narrow way between houses; \kn{ಬೀದಿ} is a proper street.

\begin{cousinweb}
\textbf{\kn{ಓಣಿ}} (\emph{ōṇi}) is native Dravidian: DEDR 1046 lists Kannada \emph{ōṇi}, \textquotedblleft{}a lane, an alley, a street\textquotedblright{}.
\end{cousinweb}

\begin{grammarlens}[title={What you've built}]
One more word for getting around.
\end{grammarlens}

\subsection*{Your turn}
\begin{itemize}
\raggedright
  \item \emph{Say it:} \emph{ōṇi}
  \item \emph{Say it:} \emph{ōṇi}, once more
  \item \emph{Say it:} \emph{haḷeya ūrina ōṇigaḷu tumbā kiridu.} — The lanes of the old town are very narrow.
  \item \emph{Recall:} say the Kannada for a landmark, then the Kannada for to view, then say \emph{ōṇi} again
\end{itemize}

\begin{tcolorbox}[breakable,colback=teal!4,colframe=teal!35!black,title={Before you move on}]
What does \kn{ಓಣಿ} mean? (\textquotedblleft{}A lane; an alley\textquotedblright{}.) Say it once more.
\end{tcolorbox}
\section[\texorpdfstring{\kn{ಹಾದುಹೋಗು} (hāduhōgu)}{ಹಾದುಹೋಗು (hāduhōgu)}]{\kn{ಹಾದುಹೋಗು} (hāduhōgu) — to pass through; to pass by}

\label{lesson:KA-C457-haduhogu}



\begin{quote}
Before the new one: say the Kannada for to view, then the Kannada for a lane.
\end{quote}

\subsection*{You'll want to know: \kn{ಹಾದುಹೋಗು}}
\textbf{\kn{ಹಾದುಹೋಗು}} — \emph{hāduhōgu} — \textquotedblleft{}to pass through; to pass by\textquotedblright{}.

\begin{quote}
\kn{ನಮ್ಮ} \kn{ರೈಲು} \kn{ಸುರಂಗದ} \kn{ಮೂಲಕ} \kn{ಹಾದುಹೋಯಿತು}.
\end{quote}

\begin{quote}
\emph{namma railu suraṅgada mūlaka hāduhōyitu.}
\end{quote}

\begin{quote}
Our train passed through a tunnel.
\end{quote}

Past (I …): \kn{ಹಾದುಹೋದೆ}, \emph{hāduhōde}.

\kn{ಮೂಲಕ} \kn{ಹಾದುಹೋಗು}, to pass through; \kn{ಪಕ್ಕದಿಂದ} \kn{ಹಾದುಹೋಗು}, to pass by.

\begin{cousinweb}
\textbf{\kn{ಹಾದುಹೋಗು}} (\emph{hāduhōgu}) is native Dravidian: \emph{hādu}, the participle of \emph{hāy}, \textquotedblleft{}to cross, to step over\textquotedblright{} (DEDR 4087), with \kn{ಹೋಗು}, to go.
\end{cousinweb}

\begin{grammarlens}[title={What you've built}]
One more word for getting around.
\end{grammarlens}

\subsection*{Your turn}
\begin{itemize}
\raggedright
  \item \emph{Say it:} \emph{hāduhōgu}
  \item \emph{Say it:} \emph{hāduhōgu}, once more
  \item \emph{Say it:} \emph{namma railu suraṅgada mūlaka hāduhōyitu.} — Our train passed through a tunnel.
  \item \emph{Recall:} say the Kannada for to view, then the Kannada for a lane, then say \emph{hāduhōgu} again
\end{itemize}

\begin{tcolorbox}[breakable,colback=teal!4,colframe=teal!35!black,title={Before you move on}]
What does \kn{ಹಾದುಹೋಗು} mean? (\textquotedblleft{}To pass through; to pass by\textquotedblright{}.) Say it once more.
\end{tcolorbox}
\section[\texorpdfstring{\kn{ಅನ್ವೇಷಿಸು} (anvēṣisu)}{ಅನ್ವೇಷಿಸು (anvēṣisu)}]{\kn{ಅನ್ವೇಷಿಸು} (anvēṣisu) — to explore; to search out}

\label{lesson:KA-C457-anvesisu}



\begin{quote}
Before the new one: say the Kannada for a lane, then the Kannada for to pass through.
\end{quote}

\subsection*{You'll want to know: \kn{ಅನ್ವೇಷಿಸು}}
\textbf{\kn{ಅನ್ವೇಷಿಸು}} — \emph{anvēṣisu} — \textquotedblleft{}to explore; to search out\textquotedblright{}.

\begin{quote}
\kn{ನಾವು} \kn{ಕಾಡಿನ} \kn{ಹೊಸ} \kn{ದಾರಿಗಳನ್ನು} \kn{ಅನ್ವೇಷಿಸಿದೆವು}.
\end{quote}

\begin{quote}
\emph{nāvu kāḍina hosa dārigaḷannu anvēṣisidevu.}
\end{quote}

\begin{quote}
We explored new paths in the forest.
\end{quote}

Past (I …): \kn{ಅನ್ವೇಷಿಸಿದೆ}, \emph{anvēṣiside}.

What is explored takes -\kn{ಅನ್ನು}; \kn{ಅನ್ವೇಷಣೆ} is exploration.

\begin{cousinweb}
\textbf{\kn{ಅನ್ವೇಷಿಸು}} (\emph{anvēṣisu}) is built from Sanskrit \emph{anvēṣaṇa}, \textquotedblleft{}seeking, searching for\textquotedblright{} (Monier-Williams), with the verb ending -\emph{isu}.
\end{cousinweb}

\begin{grammarlens}[title={What you've built}]
One more word for getting around.
\end{grammarlens}

\subsection*{Your turn}
\begin{itemize}
\raggedright
  \item \emph{Say it:} \emph{anvēṣisu}
  \item \emph{Say it:} \emph{anvēṣisu}, once more
  \item \emph{Say it:} \emph{nāvu kāḍina hosa dārigaḷannu anvēṣisidevu.} — We explored new paths in the forest.
  \item \emph{Recall:} say the Kannada for a lane, then the Kannada for to pass through, then say \emph{anvēṣisu} again
\end{itemize}

\begin{tcolorbox}[breakable,colback=teal!4,colframe=teal!35!black,title={Before you move on}]
What does \kn{ಅನ್ವೇಷಿಸು} mean? (\textquotedblleft{}To explore; to search out\textquotedblright{}.) Say it once more.
\end{tcolorbox}
\section[\texorpdfstring{\kn{ತಾಣ} (tāṇa)}{ತಾಣ (tāṇa)}]{\kn{ತಾಣ} (tāṇa) — a spot; a haunt}

\label{lesson:KA-C457-tana}



\begin{quote}
Before the new one: say the Kannada for to pass through, then the Kannada for to explore.
\end{quote}

\subsection*{You'll want to know: \kn{ತಾಣ}}
\textbf{\kn{ತಾಣ}} — \emph{tāṇa} — \textquotedblleft{}a spot; a haunt\textquotedblright{}.

\begin{quote}
\kn{ಈ} \kn{ಜಲಪಾತ} \kn{ಪ್ರವಾಸಿಗರ} \kn{ಅಚ್ಚುಮೆಚ್ಚಿನ} \kn{ತಾಣ}.
\end{quote}

\begin{quote}
\emph{ī jalapāta pravāsigara accumeccina tāṇa.}
\end{quote}

\begin{quote}
This waterfall is a favourite spot for tourists.
\end{quote}

\kn{ಪ್ರವಾಸಿ} \kn{ತಾಣ}, a tourist spot.

\begin{cousinweb}
\textbf{\kn{ತಾಣ}} (\emph{tāṇa}) is usually traced to Sanskrit \emph{sthāna}, \textquotedblleft{}a place\textquotedblright{}, through a Prakrit form \emph{ṭhāṇa} (CDIAL 13753).
\end{cousinweb}

\begin{grammarlens}[title={What you've built}]
One more word for getting around.
\end{grammarlens}

\subsection*{Your turn}
\begin{itemize}
\raggedright
  \item \emph{Say it:} \emph{tāṇa}
  \item \emph{Say it:} \emph{tāṇa}, once more
  \item \emph{Say it:} \emph{ī jalapāta pravāsigara accumeccina tāṇa.} — This waterfall is a favourite spot for tourists.
  \item \emph{Recall:} say the Kannada for to pass through, then the Kannada for to explore, then say \emph{tāṇa} again
\end{itemize}

\begin{tcolorbox}[breakable,colback=teal!4,colframe=teal!35!black,title={Before you move on}]
What does \kn{ತಾಣ} mean? (\textquotedblleft{}A spot; a haunt\textquotedblright{}.) Say it once more.
\end{tcolorbox}
\section[\texorpdfstring{\kn{ಕಂಗೊಳಿಸು} (kaṅgoḷisu)}{ಕಂಗೊಳಿಸು (kaṅgoḷisu)}]{\kn{ಕಂಗೊಳಿಸು} (kaṅgoḷisu) — to look splendid; to be a fine sight}

\label{lesson:KA-C457-kangolisu}



\begin{quote}
Before the new one: say the Kannada for to explore, then the Kannada for a spot.
\end{quote}

\subsection*{You'll want to know: \kn{ಕಂಗೊಳಿಸು}}
\textbf{\kn{ಕಂಗೊಳಿಸು}} — \emph{kaṅgoḷisu} — \textquotedblleft{}to look splendid; to be a fine sight\textquotedblright{}.

\begin{quote}
\kn{ಮಳೆಯ} \kn{ನಂತರ} \kn{ಬೆಟ್ಟಗಳು} \kn{ಹಸಿರಾಗಿ} \kn{ಕಂಗೊಳಿಸುತ್ತಿದ್ದವು}.
\end{quote}

\begin{quote}
\emph{maḷeya nantara beṭṭagaḷu hasirāgi kaṅgoḷisuttiddavu.}
\end{quote}

\begin{quote}
After the rain the hills looked splendid in green.
\end{quote}

Past (it …): \kn{ಕಂಗೊಳಿಸಿತು}, \emph{kaṅgoḷisitu}.

Said of landscapes, lights and festivals: \kn{ಊರು} \kn{ದೀಪಗಳಿಂದ} \kn{ಕಂಗೊಳಿಸಿತು}.

\begin{cousinweb}
\textbf{\kn{ಕಂಗೊಳಿಸು}} (\emph{kaṅgoḷisu}) is a Kannada verb whose origin is not settled; DEDR shows no clear source for it.
\end{cousinweb}

\begin{grammarlens}[title={What you've built}]
One more word for getting around.
\end{grammarlens}

\subsection*{Your turn}
\begin{itemize}
\raggedright
  \item \emph{Say it:} \emph{kaṅgoḷisu}
  \item \emph{Say it:} \emph{kaṅgoḷisu}, once more
  \item \emph{Say it:} \emph{maḷeya nantara beṭṭagaḷu hasirāgi kaṅgoḷisuttiddavu.} — After the rain the hills looked splendid in green.
  \item \emph{Recall:} say the Kannada for to explore, then the Kannada for a spot, then say \emph{kaṅgoḷisu} again
\end{itemize}

\begin{tcolorbox}[breakable,colback=teal!4,colframe=teal!35!black,title={Before you move on}]
What does \kn{ಕಂಗೊಳಿಸು} mean? (\textquotedblleft{}To look splendid; to be a fine sight\textquotedblright{}.) Say it once more.
\end{tcolorbox}
